% 25
\begin{frame}{Control Barrier Functions: funkcje barierowe bezpieczeństwa}
\begin{block}{Control Barrier Functions (CBF)}
Funkcje barierowe bezpieczeństwa opisują zbiór stanów dopuszczalnych i narzucają warunek, aby sterowanie nie wypchnęło układu poza ten zbiór.
\end{block}
\begin{columns}[T,onlytextwidth]
\column{0.48\textwidth}
\[
\mathcal C = \{x: h(x)\ge 0\}
\]
\[
\text{chcemy utrzymać } x(t)\in\mathcal C
\]
\column{0.48\textwidth}
\begin{itemize}
  \item $h(x)>0$: zapas bezpieczeństwa.
  \item $h(x)=0$: granica bezpieczeństwa.
  \item $h(x)<0$: stan narusza warunek.
\end{itemize}
\end{columns}
\begin{block}{Interpretacja dla nogi}
CBF nie musi uczyć chodu. Ma zabronić sterowania, które prowadzi w konfiguracje zbyt szerokie, osobliwe albo bez zapasu momentu.
\end{block}
\note{Należy użyć pełnej nazwy przed skrótem. Warto powiedzieć, że CBF to język do opisu bezpieczeństwa, a nie magiczna gwarancja bez dobrze zdefiniowanego modelu i zbioru bezpiecznego.}
\end{frame}

% 26
\begin{frame}{CBF jako filtr dla polityki uczonej}
\begin{columns}[T,onlytextwidth]
\column{0.54\textwidth}
\begin{center}
\begin{tikzpicture}[node distance=0.65cm,>=Latex]
  \tikzstyle{box}=[draw,rounded corners=2pt,fill=blue!6,minimum width=3.1cm,minimum height=0.75cm,align=center,font=\small]
  \node[box] (rl) {polityka Reinforcement\\Learning};
  \node[box,below=of rl] (filter) {filtr CBF / MPC};
  \node[box,below=of filter] (drive) {sterowanie napędów};
  \node[box,below=of drive] (robot) {robot + podłoże};
  \draw[->,thick] (rl)--node[right,font=\scriptsize]{propozycja} (filter);
  \draw[->,thick] (filter)--node[right,font=\scriptsize]{bezpieczna korekta} (drive);
  \draw[->,thick] (drive)--(robot);
  \draw[->,thick] (robot.west) .. controls +(-1.2,0) and +(-1.2,0) .. (filter.west);
\end{tikzpicture}
\end{center}
\column{0.42\textwidth}
\begin{itemize}
  \item RL proponuje szybkie zachowanie.
  \item Filtr sprawdza warunki bezpieczeństwa.
  \item Jeśli akcja jest ryzykowna, filtr ją koryguje lub odrzuca.
\end{itemize}
\begin{block}{Najbliższa literatura}
CBF + MPC zademonstrowano m.in. na ANYmalu w pracy o wielowarstwowym bezpieczeństwie lokomocji.
\end{block}
\end{columns}
\note{To slajd z bardzo praktycznym przekazem: nie trzeba od razu zastępować RL. Można zacząć od nadzorcy, który ma jawny model i jawne bariery.}
\end{frame}

% 27
\begin{frame}{Docelowa architektura projektu}
\begin{center}
\begin{tikzpicture}[node distance=0.45cm and 0.42cm,>=Latex]
  \tikzstyle{box}=[draw,rounded corners=2pt,fill=blue!6,minimum width=1.95cm,minimum height=0.62cm,align=center,font=\scriptsize]
  \node[box] (sensors) {czujniki};
  \node[box,right=of sensors] (est) {estymacja\\stanu};
  \node[box,right=of est] (rl) {polityka\\RL};
  \node[box,right=of rl] (mpc) {MPC / filtr\\bezpieczeństwa};
  \node[box,right=of mpc] (drive) {sterowanie\\impedancyjne};
  \node[box,right=of drive] (robot) {robot};
  \node[box,below=of mpc] (cbf) {warstwa CBF};
  \draw[->,thick] (sensors)--(est);
  \draw[->,thick] (est)--(rl);
  \draw[->,thick] (rl)--(mpc);
  \draw[->,thick] (cbf)--(mpc);
  \draw[->,thick] (mpc)--(drive);
  \draw[->,thick] (drive)--(robot);
  \draw[->,thick] (est.south) |- (cbf.west);
  \draw[->,thick] (robot.south) .. controls +(0,-0.9) and +(0,-0.9) .. (sensors.south);
\end{tikzpicture}
\end{center}
\vspace{0.2em}
\begin{itemize}
  \item Wersja minimalna: CBF/MPC filtruje ryzykowne konfiguracje jednej nogi.
  \item Wersja rozszerzona: MPC planuje siły kontaktowe, a WBC rozdziela zadania na przeguby.
  \item Warunek projektu naukowego: mierzalna redukcja awarii poza scenariuszami treningowymi.
\end{itemize}
\note{To jest slajd spinający projekt. Można na nim powiedzieć: nie budujemy wszystkiego naraz. Pierwszy etap to mały filtr bezpieczeństwa i model jednej nogi, dopiero później pełna lokomocja.}
\end{frame}

% 28
\begin{frame}{Demonstracja czy wynik badawczy?}
\begin{columns}[T,onlytextwidth]
\column{0.48\textwidth}
\begin{block}{Demonstracja}
\begin{itemize}
  \item Robot przeszedł trasę.
  \item Film wygląda przekonująco.
  \item Sukces może zależeć od jednego podłoża, jednej baterii i jednego strojenia.
\end{itemize}
\end{block}
\column{0.48\textwidth}
\begin{block}{Wynik badawczy}
\begin{itemize}
  \item Zdefiniowano klasę awarii.
  \item Metoda mierzalnie redukuje tę klasę awarii.
  \item Testy obejmują podłoża, zakłócenia i przypadki spoza treningu.
\end{itemize}
\end{block}
\end{columns}
\vspace{0.4em}
\begin{block}{Proponowane kryterium pierwszego etapu}
Na jednej nodze: mniej wejść w konfiguracje nieodzyskiwalne przy ograniczeniach tarcia, momentu i geometrii, bez utraty zdolności do osiągania punktów roboczych.
\end{block}
\note{Użytkownik chciał mocniej podkreślić różnicę między projektem demonstracyjnym i naukowym. Tu jest sedno: film nie wystarcza. Trzeba mieć definicję awarii, metrykę i porównanie z bazą.}
\end{frame}

% 29 Appendix A
\appendix
\begin{frame}{Appendix A: od sił kontaktowych do momentów napędów}
\begin{columns}[T,onlytextwidth]
\column{0.48\textwidth}
\begin{block}{Najprostsza intuicja statyczna}
\[
\tau \approx J(q)^T f
\]
\end{block}
\begin{itemize}
  \item $f$ - siła działająca w stopie.
  \item $J(q)^T$ - przeniesienie siły na momenty przegubów.
\end{itemize}
\column{0.48\textwidth}
\begin{block}{W rzeczywistości dochodzi}
\begin{itemize}
  \item dynamika członów,
  \item więzy mechanizmu równoległego,
  \item tarcie i poślizg,
  \item limity napędów i przekładni,
  \item opóźnienia komunikacji i regulatorów.
\end{itemize}
\end{block}
\end{columns}
\begin{block}{Dlaczego appendix?}
Ten slajd jest dobry do dyskusji technicznej, ale nie jest potrzebny każdemu odbiorcy w głównej narracji.
\end{block}
\note{Appendix pomaga, jeśli ktoś zapyta: skąd momenty w silnikach biorą się z planowania sił kontaktowych. To też miejsce na dyskusję, czy model jednej nogi wystarczy.}
\end{frame}

% 30 Appendix B
\begin{frame}{Appendix B: materiały startowe i pełniejszy problem}
\small
\begin{itemize}
  \item \href{https://underactuated.mit.edu/trajopt.html}{MIT Underactuated Robotics: Trajectory Optimization} - najlepszy pierwszy materiał o optymalizacji trajektorii i idei przesuwającego się horyzontu.
  \item \href{https://ieeexplore.ieee.org/document/8594448}{MIT Cheetah 3 przez wypukłe MPC} - klasyczny przykład szybkiego planowania sił reakcji podłoża dla czworonoga.
  \item \href{https://github.com/iit-DLSLab/Quadruped-PyMPC}{Quadruped-PyMPC} - praktyczny kod w Pythonie dla MPC opartego o model pojedynczego ciała sztywnego.
  \item \href{https://arxiv.org/abs/2011.00032}{CBF + MPC dla ANYmal} - najbliższy punkt startowy dla warstwy bezpieczeństwa w robotach kroczących.
  \item \href{https://arxiv.org/abs/2503.04613}{Whole-Body MPC z MuJoCo} - współczesny przykład MPC całego ciała dla czworonogów i humanoidów.
\end{itemize}
\vspace{-0.25em}
\begin{block}{Pełniejsza formulacja jednej nogi}
\[
\min \sum_k \left(
\|p_F-p_F^{ref}\|^2 + \|q-q_{nom}\|^2 + \|\tau\|^2\right)
\]
\[
\Phi(q,z)=0,\quad q_{min}\le q\le q_{max},\quad |\tau|\le\tau_{max},\quad h(q,\dot q)\ge0.
\]
\end{block}
\note{To jest slajd kończący i jednocześnie źródła. Każdy materiał ma jedno zdanie opisu, tak jak użytkownik poprosił. Dyad jest świadomie pominięty ze względu na komercyjną licencję, a może być wspomniany ustnie.}
\end{frame}
